% =====================================================================
%  Affinity-map / first-pass synthesis BODY (Field Day 1, 2026-06-19).
%  Requires affinity-map-preamble.tex loaded in the parent's preamble.
%  In the standalone handout it renders with a title block; when embedded
%  (main_id.tex defines \affinityembedded) the big title is suppressed so
%  the appendix \chapter heading provides the title instead.
% =====================================================================
\setlist[itemize]{leftmargin=1.2em, itemsep=1pt, topsep=2pt}

\ifdefined\affinityembedded\else
\begin{center}
  {\footnotesize\textsc{Urban Innovation Challenge \textbullet{} CBL Investigate $\rightarrow$ Refine \textbullet{} Field Day 1 (Thu 18 Jun 2026)}}\\[2pt]
  {\LARGE\bfseries Walking the Sudirman Corridor}\\[1pt]
  {\large Affinity Map \& First-Pass Synthesis}
\end{center}
\vspace{2pt}
\fi

\begin{tcolorbox}[colback=ink,colframe=ink,arc=2pt,boxsep=3pt,left=8pt,right=8pt,top=5pt,bottom=5pt]
  \textcolor{white}{\textbf{\textcolor{yellow!80!orange}{Refined challenge (proposed):}}
  Make walking the Sudirman corridor \textbf{comfortable and safe} --- through the heat of the
  day, into the night, and beyond the main trotoar to the streets one block behind it.}
\end{tcolorbox}

{\footnotesize\color{ink!70}16 people engaged (10 upper/middle \textbullet{} 6 worker/``bottom'') \textbullet{}
12 recorded intercepts + notulen \textbullet{} 105 geotagged photos/videos \textbullet{} 6 weather readings
\textbullet{} daytime only (10:00--15:00). Scores below are \textbf{analyst-interpreted from transcripts},
not measured.\par}

% ---------------- 1. WE ENGAGED ----------------
\sect{1\quad We engaged}
\begin{tcbraster}[raster columns=3, raster equal height=rows, colback=white, colframe=soft,
  fonttitle=\bfseries\small, coltitle=ink, size=small, arc=2pt]
  \begin{tcolorbox}[title=Pedestrians]\footnotesize
    \begin{itemize}\item Daily \& weekly commuters (KRL/MRT/TJ)\item Elderly daily walker (1.5\,h exercise)
    \item Parent + child; students (SMAN 3)\item Motorbike rider (walks last hop only)\end{itemize}
  \end{tcolorbox}
  \begin{tcolorbox}[title=Workers of the corridor]\footnotesize
    \begin{itemize}\item Area \& JPO security guards\item MITJ staff (incl.\ tunnel post)
    \item Satpol PP; a PKL (Starling)\item GrabCar driver; local elder (Kebon Melati)\end{itemize}
  \end{tcolorbox}
  \begin{tcolorbox}[title=Where / when]\footnotesize
    \begin{itemize}\item Main trotoar, stations, JPO, Kendal tunnel\item GPS: Dukuh Atas $\to$ Setiabudi/Semanggi
    \item \textbf{Gap:} back-streets barely sampled\item \textbf{Gap:} no night, no rain\end{itemize}
  \end{tcolorbox}
\end{tcbraster}

% ---------------- 2. WE FOUND ----------------
\sect{2\quad We found --- affinity clusters}
\begin{tcbraster}[raster columns=2, raster equal height=rows, colback=white, colframe=soft,
  fonttitle=\bfseries\small, coltitle=ink, size=small, arc=2pt]
  \begin{tcolorbox}[title={A good sidewalk, undermined by heat\tg{neg}{PAIN \textbullet{} dominant}}]\footnotesize
    \textit{``Trotoarnya udah bagus\dots udah manusiawi? Udah.''} --- yet ``panas'' is the \#1 complaint.
    \begin{itemize}\item Main trotoar widely praised: \textit{lega, manusiawi, bagus}.
    \item Heat named by \textbf{7 of 12}; design critiqued as copied from 4-season countries.
    \item Coping: shade-seeking, umbrellas, cooler hours.\end{itemize}
  \end{tcolorbox}
  \begin{tcolorbox}[title={Safe by day, risky at isolated night spots\tg{mix}{PAIN \textbullet{} conditional}}]\footnotesize
    \begin{itemize}\item Daytime feels safe (incl.\ a woman walking home 01:00).
    \item Hotspots: JPO \& parking after $\sim$22:00 --- \textbf{lit but empty}, not dark.
    \item Guard: \textit{``terang\dots cuma sepi.''} Crime down $\sim$6 months.\end{itemize}
  \end{tcolorbox}
  \begin{tcolorbox}[title={Micro-hazards \& conflicts\tg{neg}{PAIN}}]\footnotesize
    \begin{itemize}\item Kerb$\to$zebra ramp too steep/slippery --- a respondent's friend \textbf{fell twice}.
    \item Ojol/Gojek riding onto the trotoar (near-misses).
    \item PKL (Starling) --- convenience vs obstruction tension.\end{itemize}
  \end{tcolorbox}
  \begin{tcolorbox}[title={Institutional gaps\tg{neg}{PAIN \textbullet{} systemic}}]\footnotesize
    \begin{itemize}\item Security responsibility falls between KAI/MRT/police zones.
    \item Broken/uneven infra reported but \textit{``gak pernah digubris.''}
    \item Open Q: is JAKI reporting working? UNHCR donation scam seen.\end{itemize}
  \end{tcolorbox}
\end{tcbraster}

% ---------------- 3. QUANTITATIVE ----------------
\sect{3\quad Quantitative --- what the numbers show}
\noindent
\begin{minipage}[t]{0.49\textwidth}
  \begin{center}\footnotesize\textbf{Measured weather across the session}\\[-2pt]
  \begin{tikzpicture}
    \begin{axis}[width=\linewidth,height=4.6cm,ymin=30,ymax=36,xmin=0.6,xmax=6.4,
      xtick={1,2,3,4,5,6},xticklabels={10:20,10:30,10:55,12:48,13:03,13:35},
      x tick label style={font=\scriptsize,rotate=30,anchor=east},
      ylabel={\scriptsize Temp ($^{\circ}$C)},ytick={30,32,34,36},
      tick label style={font=\scriptsize},grid=major,grid style={soft},
      legend style={font=\scriptsize,at={(0.5,1.18)},anchor=south,legend columns=2,draw=none}]
      \addplot[heat,very thick,mark=*] coordinates {(1,33)(2,34)(3,34)(4,35)(5,35)(6,35)};
      \addplot[accent,very thick,dashed,mark=square*] coordinates {(1,31)(2,31)(3,32)(4,33)(5,33)(6,33)};
      \legend{Feels-like,Air temp}
    \end{axis}
  \end{tikzpicture}\\[-1pt]
  \scriptsize UV 7--8 (Very High) \textbullet{} 0\,mm rain \textbullet{} +3$^{\circ}$C above avg high.
  Feels-like held \textbf{above the $\sim$27$^{\circ}$C comfort threshold all day}.
  \end{center}
\end{minipage}\hfill
\begin{minipage}[t]{0.49\textwidth}
  \begin{center}\footnotesize\textbf{Themes mentioned (of 12 recorded interviews)}\end{center}
  \footnotesize
  \begin{tabular}{@{}l l r@{}}
    Heat / too hot        & \sbar{heat}{4.06cm} & 7\\
    Sidewalk is good/wide & \sbar{pos}{3.48cm}  & 6\\
    Feels safe (daytime)  & \sbar{pos}{3.48cm}  & 6\\
    Night/isolation risk  & \sbar{neg}{2.32cm}  & 4\\
    Crossing ramp hazard  & \sbar{neg}{1.16cm}  & 2\\
    Institutional gap     & \sbar{neg}{1.16cm}  & 2\\
    Ojol on trotoar       & \sbar{neg}{0.58cm}  & 1\\
  \end{tabular}\\[2pt]
  {\scriptsize Qualitative coding of transcripts; a person can raise several themes.}
\end{minipage}

\vspace{6pt}
\noindent
\begin{minipage}[t]{0.46\textwidth}
  \footnotesize\textbf{Interpreted self-rated scores by group (1--10)}\\[2pt]
  \begin{tabular}{@{}lccc@{}}
    \toprule
    Group (n) & Comfort & Security & Heat-disc.\\
    \midrule
    Upper/middle (8) & 7.0 & 8.1 & 7.4\\
    Worker (4)       & 7.5 & 6.3 & 5.8\\
    All (12)         & 7.2 & 7.5 & 6.8\\
    \bottomrule
  \end{tabular}\\[2pt]
  {\scriptsize Heat-disc.: 1=cool, 10=hot.}
\end{minipage}\hfill
\begin{minipage}[t]{0.51\textwidth}
  \begin{tcolorbox}[colback=yellow!8,colframe=mix!60,arc=2pt,boxsep=2pt,left=6pt,right=6pt,top=4pt,bottom=4pt]
    \footnotesize\textbf{\textcolor{mix!70!black}{$\triangle$ Read as hypotheses, not results.}}
    Scores are \textbf{interpreted from transcripts, not asked}, n is tiny. The apparent pattern
    (workers rate security \& heat lower) is a \textbf{lead for Saturday's measured scales}.
    A Mann--Whitney / permutation test on the interpreted heat scores gives effect size
    $r=-0.56$ but $p\approx0.15$ --- \textbf{not significant}; \textbf{no t-test is reported}
    (see \texttt{analysis-plan.md}).
  \end{tcolorbox}
\end{minipage}

% ---------------- 4. SENTIMENT ----------------
\sect{4\quad Sentiment by theme (transparent hand-coding)}
\noindent\footnotesize
\begin{tabular}{@{}l l l @{\hspace{2.2em}} l l l@{}}
  Sidewalk / infrastructure & \sbar{pos}{2.6cm} & \textcolor{pos}{$+$} &
  Heat / thermal comfort    & \sbar{neg}{2.6cm} & \textcolor{neg}{$-$}\\
  Greenery / shade buffer   & \sbar{pos}{2.3cm} & \textcolor{pos}{$+$} &
  Safety --- night/isolated & \sbar{neg}{2.1cm} & \textcolor{neg}{$-$}\\
  Safety --- daytime        & \sbar{pos}{2.5cm} & \textcolor{pos}{$+$} &
  Maintenance / complaints  & \sbar{neg}{1.8cm} & \textcolor{neg}{$-$}\\
  & & & Enforcement / ojol   & \sbar{neg}{1.3cm} & \textcolor{neg}{$-$}\\
\end{tabular}\\[2pt]
{\scriptsize Direction of sentiment per theme, coded by hand (not an automated score). The split:
\textbf{positive on the built sidewalk, negative on heat, night-isolation, and upkeep}.}

% ---------------- 5. OPEN QUESTIONS ----------------
\sect{5\quad Open questions --- what Saturday must test}
\begin{tcbraster}[raster columns=4, raster equal height=rows, colback=accent!5, colframe=accent!40,
  fonttitle=\bfseries\footnotesize, coltitle=accent!70!black, size=small, arc=2pt]
  \begin{tcolorbox}[title=The gradient]\scriptsize Pair each main segment with the back-street one block behind; audit both. Is the gap as sharp as assumed? (untested)\end{tcolorbox}
  \begin{tcolorbox}[title=Night \& women]\scriptsize Does daytime ``feels safe'' hold after 22:00 at the JPO/parking? Women-led intercepts (Cikini--Benhil).\end{tcolorbox}
  \begin{tcolorbox}[title=Measured scales]\scriptsize Ask 1--10 + ASHRAE thermal vote consistently $\to$ enables real stats. 5-min counts joined to WeatherKit.\end{tcolorbox}
  \begin{tcolorbox}[title=Rain behaviour]\scriptsize Dry yesterday; $\sim$40--50\% rain forecast Sat --- observe sheltering / mode-shift.\end{tcolorbox}
\end{tcbraster}

% ---------------- FOOTER ----------------
\vspace{3pt}
\begin{tcolorbox}[colback=ink!3,colframe=soft,arc=2pt,boxsep=2pt,left=6pt,right=6pt,top=4pt,bottom=4pt]
\scriptsize\textbf{Method (data-science breakdown).}
\textbf{Qualitative} --- verbatim transcripts thematically coded into pain-points/goals (\S2).
\textbf{Quantitative} --- measured weather (\S3, real); frequencies (\S3); interpreted scales shown descriptively only.
\textbf{Sentiment} --- hand-coded direction per theme (\S4), transparent, not a black-box model.
\textbf{Inferential stats deferred:} Mann--Whitney / permutation (by group) and a comfort-vs-heat-index
correlation run on \textbf{Saturday's measured} data, not on interpreted scores. Plan: \texttt{investigate/analysis-plan.md};
evidence: \texttt{investigate/synthesis.md}. \hfill Tim Riset Pijak \textbullet{} 2026-06-19.
\end{tcolorbox}
